% !TEX program = pdflatex
% FRE6883 Recitation 2 -- Structured Programming and CRR Code Lab
\documentclass[10pt,aspectratio=169]{beamer}

% Language & encoding
\usepackage[utf8]{inputenc}
\usepackage[T1]{fontenc}
\usepackage[english]{babel}
\usepackage{lmodern}
\usepackage{microtype}
\DisableLigatures{encoding=*,family=tt*}

% Math / tables / layout
\usepackage{amsmath,amssymb,bm,mathtools}
\usepackage{booktabs}
\usepackage{multicol}
\usepackage{changepage}

% Graphics
\usepackage{graphicx}
\usepackage{tikz}
\usetikzlibrary{arrows.meta,positioning,calc,decorations.pathreplacing}
\usepackage{pgfplots}
\pgfplotsset{compat=1.18}

% Theme (Metropolis)
\usetheme[numbering=fraction,progressbar=frametitle]{metropolis}
\setbeamertemplate{frame numbering}[fraction]
\setbeamertemplate{navigation symbols}{}

% Bootcamp palette
\definecolor{BootcampBlueDark}{HTML}{1A3C68}
\definecolor{BootcampBlue}{RGB}{31,110,178}
\definecolor{AccentGold}{HTML}{DAA520}
\definecolor{SoftGray}{RGB}{245,246,248}
\definecolor{TextBlack}{HTML}{000000}
\definecolor{Crimson}{HTML}{990000}
\definecolor{MatrixGreen}{HTML}{228B22}
\definecolor{MatrixRed}{HTML}{B22222}

% Global formatting
\setbeamercolor{normal text}{fg=TextBlack,bg=white}
\setbeamercolor{title}{fg=TextBlack}
\setbeamercolor{structure}{fg=BootcampBlueDark}
\setbeamercolor{progress bar}{fg=BootcampBlueDark,bg=gray!30}

\setbeamercolor{frametitle}{bg=BootcampBlueDark,fg=white}
\setbeamerfont{frametitle}{series=\bfseries,size=\large}
\setbeamertemplate{frametitle}{
  \nointerlineskip
  \begin{beamercolorbox}[wd=\paperwidth,ht=3.2ex,dp=1.4ex,leftskip=1em,rightskip=1em]{frametitle}
    \usebeamerfont{frametitle}\insertframetitle
  \end{beamercolorbox}
}

\setbeamercolor{block title example}{bg=BootcampBlueDark!90!black,fg=white}
\setbeamercolor{block body example}{bg=BootcampBlueDark!5!white,fg=black}
\setbeamercolor{block title proposition}{bg=MatrixGreen!75!black,fg=white}
\setbeamercolor{block body proposition}{bg=MatrixGreen!7!white,fg=black}
\setbeamercolor{block title illustration}{bg=AccentGold!85!black,fg=white}
\setbeamercolor{block body illustration}{bg=AccentGold!10!white,fg=black}
\setbeamercolor{block title proof}{bg=gray!70!black,fg=white}
\setbeamercolor{block body proof}{bg=SoftGray,fg=black}
\setbeamertemplate{blocks}[rounded][shadow=false]

\newenvironment{definitionblock}[1]{%
  \par\vspace{0.3em}%
  \begin{beamercolorbox}[
      wd=\textwidth,sep=0.35em,leftskip=0.3em,rightskip=0.6em,
      rounded=false,shadow=false
    ]{block title example}%
    \raisebox{0.15em}{\usebeamerfont{block title example}#1}%
  \end{beamercolorbox}%
  \vspace{-0.4em}%
  \begin{beamercolorbox}[
      wd=\textwidth,sep=0.55em,leftskip=0.6em,rightskip=0.6em,
      rounded=false,shadow=false
    ]{block body example}%
    \usebeamerfont{block body example}%
}{%
  \end{beamercolorbox}%
  \par\vspace{0.45em}%
}

\newenvironment{propositionblock}[1]{%
  \par\vspace{0.3em}%
  \begin{beamercolorbox}[
      wd=\textwidth,sep=0.35em,leftskip=0.3em,rightskip=0.6em,
      rounded=false,shadow=false
    ]{block title proposition}%
    \raisebox{0.15em}{\usebeamerfont{block title example}#1}%
  \end{beamercolorbox}%
  \vspace{-0.4em}%
  \begin{beamercolorbox}[
      wd=\textwidth,sep=0.55em,leftskip=0.6em,rightskip=0.6em,
      rounded=false,shadow=false
    ]{block body proposition}%
    \usebeamerfont{block body example}%
}{%
  \end{beamercolorbox}%
  \par\vspace{0.45em}%
}

\newenvironment{illustrationblock}[1]{%
  \par\vspace{0.3em}%
  \begin{beamercolorbox}[
      wd=\textwidth,sep=0.35em,leftskip=0.3em,rightskip=0.6em,
      rounded=false,shadow=false
    ]{block title illustration}%
    \raisebox{0.15em}{\usebeamerfont{block title example}#1}%
  \end{beamercolorbox}%
  \vspace{-0.4em}%
  \begin{beamercolorbox}[
      wd=\textwidth,sep=0.55em,leftskip=0.6em,rightskip=0.6em,
      rounded=false,shadow=false
    ]{block body illustration}%
    \usebeamerfont{block body example}%
}{%
  \end{beamercolorbox}%
  \par\vspace{0.45em}%
}

\newenvironment{proofblock}[1]{%
  \par\vspace{0.3em}%
  \begin{beamercolorbox}[
      wd=\textwidth,sep=0.35em,leftskip=0.3em,rightskip=0.6em,
      rounded=false,shadow=false
    ]{block title proof}%
    \raisebox{0.15em}{\usebeamerfont{block title example}#1}%
  \end{beamercolorbox}%
  \vspace{-0.4em}%
  \begin{beamercolorbox}[
      wd=\textwidth,sep=0.55em,leftskip=0.6em,rightskip=0.6em,
      rounded=false,shadow=false
    ]{block body proof}%
    \usebeamerfont{block body example}%
}{%
  \end{beamercolorbox}%
  \par\vspace{0.45em}%
}

% C++ listings follow the Lecture 2 and Recitation 1 conventions.
\usepackage{listings}
\usepackage{textcomp}
\lstdefinestyle{coursecpp}{
  language=C++,
  basicstyle=\ttfamily\fontsize{9.2}{10.5}\selectfont,
  keywordstyle=\color{blue},
  commentstyle=\color{MatrixGreen},
  stringstyle=\color{Crimson},
  showstringspaces=false,
  columns=fullflexible,
  keepspaces=true,
  tabsize=4,
  breaklines=false,
  frame=none,
  aboveskip=0.25em,
  belowskip=0.25em,
  upquote=true
}
\lstdefinestyle{solutioncpp}{
  style=coursecpp,
  basicstyle=\ttfamily\fontsize{8.3}{9.4}\selectfont
}
\lstdefinestyle{tinycpp}{
  style=coursecpp,
  basicstyle=\ttfamily\fontsize{7.7}{8.7}\selectfont
}
\lstset{style=coursecpp}
\setbeamercovered{invisible}
\setlength{\parskip}{0.1em}

\newcommand{\problemmeta}[2]{%
 {\color{BootcampBlue}\small\textbf{GUIDED CODE}\quad #1\hfill #2}\par\smallskip\small}
\newcommand{\solutioninfo}[1]{%
 {\color{MatrixGreen!75!black}\textbf{#1 -- Solution}}\par\smallskip}
\newcommand{\explain}[2]{%
 {\color{BootcampBlueDark}\textbf{#1}}\par #2\par\medskip}

\setbeamertemplate{frame footer}{\scriptsize FRE6883 | Recitation 2 | Structured Programming and CRR}
\title{Structured Programming and CRR Code Lab}
\author{}
\date{}

% 38 core frames, planned for approximately 100 minutes, plus reserve practice.
% Source baseline: Lecture 2, Topic 2, pp. 7--37.
% Visual and exercise conventions: Recitation 1, Decks 1 and 5.
% Model: gross factors, S0 > 0, 0 < D < R < U, European call.
% Technical correction: no-arbitrage rejection uses R <= D || R >= U.
% Standard C++ correction: fixed capacity replaces Lecture 2's VLA Price[N+1].
\begin{document}

% FRAME 1 -- 1 minute
\begin{frame}[plain]
\vfill
{\usebeamerfont{title}\inserttitle\par}
\medskip
{\color{BootcampBlueDark}\rule{0.95\textwidth}{0.8pt}}
\bigskip
{\Large FRE6883 -- Recitation 2\par}
\medskip
Complete it, trace it, test it, explain it
\vfill
\end{frame}

% FRAME 2 -- 2 minutes
\begin{frame}{Today we rebuild one pricer from small testable pieces}
\begin{columns}[c,onlytextwidth]
\begin{column}{0.45\textwidth}
\centering
Model data $S_0,U,D,R,N,K$\par
$\downarrow$\par
Stock-price nodes\par
$\downarrow$\par
Terminal payoffs\par
$\downarrow$\par
Backward induction\par
$\downarrow$\par
\textbf{Option price $H(0,0)$}
\end{column}
\begin{column}{0.51\textwidth}
\begin{propositionblock}{Your job on every activity}
Predict the missing operation, complete the code, trace one example, and name one test that could expose a bug.
\end{propositionblock}
\medskip
The goal is not to memorize the final function. It is to make every line explainable.
\end{column}
\end{columns}
\end{frame}

% FRAME 3 -- 3 minutes
\begin{frame}[fragile]{Three control structures are enough to express the algorithm}
\begin{columns}[T,onlytextwidth]
\begin{column}{0.31\textwidth}
\begin{definitionblock}{Sequence}
\begin{lstlisting}[basicstyle=\ttfamily\fontsize{8}{9}\selectfont]
read inputs;
compute node;
print result;
\end{lstlisting}
One statement follows another.
\end{definitionblock}
\end{column}
\begin{column}{0.31\textwidth}
\begin{definitionblock}{Condition}
\begin{lstlisting}[basicstyle=\ttfamily\fontsize{8}{9}\selectfont]
if (invalid)
{
    return -1;
}
\end{lstlisting}
The data decides the path.
\end{definitionblock}
\end{column}
\begin{column}{0.34\textwidth}
\begin{definitionblock}{Iteration}
\begin{lstlisting}[basicstyle=\ttfamily\fontsize{8}{9}\selectfont]
for (int i = 0;
     i <= n; ++i)
{
    visit(i);
}
\end{lstlisting}
The loop repeats one rule.
\end{definitionblock}
\end{column}
\end{columns}
\medskip
\textbf{Warm-up:} point to one of each structure in the CRR procedure.
\end{frame}

% FRAME 4 -- Activity 1 problem
\begin{frame}[fragile]{01 | Compute one stock-price node -- Problem}
\problemmeta{Formula to function}{Complete together}
At step $n$ and node $i$,
\[
S(n,i)=S_0U^iD^{\,n-i}, \qquad 0\leq i\leq n.
\]
\begin{lstlisting}
double CalculateAssetPrice(
    double S0, double U, double D, int n, int i)
{
    return /* complete the expression */;
}
\end{lstlisting}
\begin{definitionblock}{Hand check}
$S_0=106$, $U=1.15125$, $D=0.86862$, $n=3$, $i=2$ should give approximately $122.03$.
\end{definitionblock}
{\small\color{BootcampBlueDark}\textbf{Think first:} Which exponent counts up moves? Which counts down moves?}
\end{frame}

% FRAME 5 -- Activity 1 solution
\begin{frame}[fragile]{01 | Compute one stock-price node -- Solution}
\begin{columns}[T,onlytextwidth]
\begin{column}{0.62\textwidth}
\begin{lstlisting}[style=solutioncpp]
double CalculateAssetPrice(
    double S0, double U, double D, int n, int i)
{
    return S0 * std::pow(U, i)
              * std::pow(D, n - i);
}
\end{lstlisting}
\end{column}
\begin{column}{0.35\textwidth}
\small
\explain{The idea}{$i$ counts up moves; $n-i$ counts down moves. Their sum is always $n$.}
\explain{Boundary nodes}{$i=0$ gives $S_0D^n$; $i=n$ gives $S_0U^n$.}
\explain{Required header}{\texttt{std::pow} is declared in \texttt{<cmath>}.}
\end{column}
\end{columns}
\end{frame}

% FRAME 6 -- Activity 2 problem
\begin{frame}[fragile]{02 | Visit every node exactly once -- Problem}
\problemmeta{Nested loops}{Complete the bounds}
For $N=3$, the valid node pairs are
\[
(0,0);\ (1,0),(1,1);\ (2,0),(2,1),(2,2);\ (3,0),(3,1),(3,2),(3,3).
\]
\begin{lstlisting}
for (int n = /* start */; n /* condition */ N; ++n)
{
    for (int i = /* start */; i /* condition */ n; ++i)
    {
        std::cout << CalculateAssetPrice(
            S0, U, D, n, i) << std::endl;
    }
}
\end{lstlisting}
{\small\color{BootcampBlueDark}\textbf{Think first:} How many nodes should the code print when $N=3$?}
\end{frame}

% FRAME 7 -- Activity 2 solution
\begin{frame}[fragile]{02 | Visit every node exactly once -- Solution}
\begin{columns}[T,onlytextwidth]
\begin{column}{0.58\textwidth}
\begin{lstlisting}[style=solutioncpp]
for (int n = 0; n <= N; ++n)
{
    for (int i = 0; i <= n; ++i)
    {
        std::cout << CalculateAssetPrice(
            S0, U, D, n, i) << std::endl;
    }
}
\end{lstlisting}
\end{column}
\begin{column}{0.38\textwidth}
\small
\explain{Outer loop}{Selects one time step from $0$ through $N$.}
\explain{Inner loop}{Visits its $n+1$ valid nodes, from $0$ through $n$.}
\explain{Node count}{$1+2+3+4=10$ when $N=3$; in general $(N+1)(N+2)/2$.}
\end{column}
\end{columns}
\end{frame}

% FRAME 8 -- Activity 3 problem
\begin{frame}[fragile]{03 | Store the visited nodes in one array -- Problem}
\problemmeta{Array + running index}{Complete three blanks}
\begin{lstlisting}[style=solutioncpp]
const int MAX_N = 8;
const int SIZE = (MAX_N + 1) * (MAX_N + 2) / 2;
double prices[SIZE] = { 0.0 };

int index = /* ? */;
for (int n = 0; n <= MAX_N; ++n)
{
    for (int i = 0; i <= n; ++i)
    {
        prices[index] = CalculateAssetPrice(
            S0, U, D, n, i);
        /* advance index */;
    }
}
\end{lstlisting}
{\small\color{BootcampBlueDark}\textbf{Think first:} What is the first valid array index? What is the last?}
\end{frame}

% FRAME 9 -- Activity 3 solution
\begin{frame}[fragile]{03 | Store the visited nodes in one array -- Solution}
\begin{columns}[T,onlytextwidth]
\begin{column}{0.60\textwidth}
\begin{lstlisting}[style=solutioncpp]
const int MAX_N = 8;
const int SIZE = (MAX_N + 1)
               * (MAX_N + 2) / 2;
double prices[SIZE] = { 0.0 };

int index = 0;
for (int n = 0; n <= MAX_N; ++n)
{
    for (int i = 0; i <= n; ++i)
    {
        prices[index] = CalculateAssetPrice(
            S0, U, D, n, i);
        ++index;
    }
}
\end{lstlisting}
\end{column}
\begin{column}{0.37\textwidth}
\small
\explain{Capacity}{Steps $0$ through $8$ contain $9\cdot10/2=45$ nodes.}
\explain{Invariant}{Before each write, \texttt{index} identifies the next free slot.}
\explain{Safety}{The final write uses index $44$; then \texttt{index} becomes $45$.}
\end{column}
\end{columns}
\end{frame}

% FRAME 10 -- checkpoint
\begin{frame}{A loop generates nodes; an array remembers them}
\begin{columns}[c,onlytextwidth]
\begin{column}{0.47\textwidth}
\begin{definitionblock}{Computation}
\texttt{CalculateAssetPrice} translates one mathematical rule into one reusable operation.
\end{definitionblock}
\begin{definitionblock}{Traversal}
The nested loops generate every valid pair $(n,i)$.
\end{definitionblock}
\end{column}
\begin{column}{0.49\textwidth}
\begin{definitionblock}{Storage}
The array preserves results after the loop moves to the next node.
\end{definitionblock}
\begin{illustrationblock}{Check the design}
Can we read \texttt{prices[SIZE]}? Why or why not?
\end{illustrationblock}
\end{column}
\end{columns}
\end{frame}

% FRAME 11 -- Activity 4 problem
\begin{frame}[fragile]{04 | Reject invalid model data -- Problem}
\problemmeta{Conditions + early return}{Complete the no-arbitrage test}
{\small Valid data must satisfy $S_0>0$ and $0<D<R<U$.\par}
\begin{lstlisting}[style=solutioncpp]
int ValidateInputData(
    const double& S0, const double& U,
    const double& D, const double& R)
{
    if (S0 <= 0.0 || D <= 0.0 || U <= D)
    {
        return -1;
    }
    if (/* reject R outside (D,U) */)
    {
        return -1;
    }
    return 0;
}
\end{lstlisting}
{\small\color{BootcampBlueDark}\textbf{Test:} should $D=0.9$, $R=0.9$, $U=1.1$ pass?}
\end{frame}

% FRAME 12 -- Activity 4 solution
\begin{frame}[fragile]{04 | Reject invalid model data -- Solution}
\begin{columns}[T,onlytextwidth]
\begin{column}{0.62\textwidth}
\begin{lstlisting}[style=solutioncpp]
int ValidateInputData(
    const double& S0, const double& U,
    const double& D, const double& R)
{
    if (S0 <= 0.0 || D <= 0.0 || U <= D)
    {
        return -1;
    }
    if (R <= D || R >= U)
    {
        return -1;
    }
    return 0;
}
\end{lstlisting}
\end{column}
\begin{column}{0.35\textwidth}
\small
\explain{Boundary matters}{Equality creates an invalid boundary; use \texttt{<=} and \texttt{>=}.}
\explain{Useful test}{With $D=R=0.9$, the function returns \texttt{-1}.}
\explain{Why early return}{Invalid inputs cannot reach the pricing calculation.}
\end{column}
\end{columns}
\end{frame}

% FRAME 13 -- function architecture bridge
\begin{frame}[fragile]{Functions give each financial rule one clear home}
\begin{columns}[T,onlytextwidth]
\begin{column}{0.48\textwidth}
\begin{definitionblock}{BinomialTreeModel.h}
\begin{lstlisting}[basicstyle=\ttfamily\fontsize{7.8}{8.8}\selectfont]
#pragma once
namespace fre {
double RiskNeutProb(
    double U, double D, double R);
double CalculateAssetPrice(
    double S0, double U, double D,
    int n, int i);
}
\end{lstlisting}
Declarations tell other files what exists.
\end{definitionblock}
\end{column}
\begin{column}{0.48\textwidth}
\begin{definitionblock}{BinomialTreeModel.cpp}
\begin{lstlisting}[basicstyle=\ttfamily\fontsize{7.8}{8.8}\selectfont]
#include "BinomialTreeModel.h"
namespace fre {
// definitions belong here
}
\end{lstlisting}
Definitions contain the executable rules.
\end{definitionblock}
\end{column}
\end{columns}
\medskip
\textbf{Design question:} which file should include \texttt{<cmath>}?
\end{frame}

% FRAME 14 -- Activity 5 problem
\begin{frame}[fragile]{05 | Implement the risk-neutral weight -- Problem}
\problemmeta{Small pure function}{Complete and hand-check}
\[
q=\frac{R-D}{U-D}.
\]
\begin{lstlisting}[style=solutioncpp]
namespace fre {

double RiskNeutProb(double U, double D, double R)
{
    return /* ? */;
}

}
\end{lstlisting}
\begin{definitionblock}{Hand check}
For $U=1.2$, $D=0.8$, $R=1.1$, the result should be $0.75$.
\end{definitionblock}
{\small\color{BootcampBlueDark}\textbf{Think first:} Which validation condition guarantees that $0<q<1$?}
\end{frame}

% FRAME 15 -- Activity 5 solution
\begin{frame}[fragile]{05 | Implement the risk-neutral weight -- Solution}
\begin{columns}[T,onlytextwidth]
\begin{column}{0.57\textwidth}
\begin{lstlisting}[style=solutioncpp]
namespace fre {

double RiskNeutProb(double U, double D, double R)
{
    return (R - D) / (U - D);
}

}
\end{lstlisting}
\end{column}
\begin{column}{0.40\textwidth}
\small
\explain{Parentheses}{They preserve the numerator and denominator from the formula.}
\explain{Test value}{$q=(1.1-0.8)/(1.2-0.8)=0.75$.}
\explain{Precondition}{$D<R<U$ places $R-D$ strictly between $0$ and $U-D$.}
\end{column}
\end{columns}
\end{frame}

% FRAME 16 -- Activity 6 problem
\begin{frame}[fragile]{06 | Predict value and reference effects -- Problem}
\problemmeta{Trace before running}{Fill the final output}
\begin{lstlisting}[style=solutioncpp]
void Foo(int a, int& b)
{
    ++a;
    ++b;
    std::cout << a << " " << b << std::endl;
}

int main()
{
    int x = 1, y = 1;
    Foo(x, y);
    std::cout << x << " " << y << std::endl;
}
\end{lstlisting}
\begin{definitionblock}{Complete the trace}
Inside \texttt{Foo}: \texttt{\_\_ \_\_}\qquad Back in \texttt{main}: \texttt{\_\_ \_\_}
\end{definitionblock}
\end{frame}

% FRAME 17 -- Activity 6 solution
\begin{frame}[fragile]{06 | Predict value and reference effects -- Solution}
\begin{columns}[c,onlytextwidth]
\begin{column}{0.46\textwidth}
\begin{propositionblock}{Observed output}
Inside \texttt{Foo}: \texttt{2 2}\par\medskip
Back in \texttt{main}: \texttt{1 2}
\end{propositionblock}
\end{column}
\begin{column}{0.50\textwidth}
\small
\explain{\texttt{int a}}{The function receives a copy of \texttt{x}. Incrementing the copy does not change \texttt{x}.}
\explain{\texttt{int\& b}}{The function receives an alias for \texttt{y}. Incrementing \texttt{b} changes \texttt{y}.}
\explain{\texttt{const int\& b}}{The alias could be read but not incremented; \texttt{++b} would fail to compile.}
\end{column}
\end{columns}
\end{frame}

% FRAME 18 -- Activity 7 problem
\begin{frame}[fragile]{07 | Return four inputs through references -- Problem}
\problemmeta{Output parameters}{Complete the reads and success code}
\begin{lstlisting}[style=solutioncpp]
int GetInputData(
    double& S0, double& U, double& D, double& R)
{
    std::cout << "Enter S0 U D R: ";
    std::cin >> /* four output parameters */;

    if (ValidateInputData(S0, U, D, R) == -1)
    {
        return /* failure */;
    }
    return /* success */;
}
\end{lstlisting}
\begin{definitionblock}{Call site}
\texttt{double s0=0.0, u=0.0, d=0.0, r=0.0;}\par
\texttt{int status = fre::GetInputData(s0, u, d, r);}
\end{definitionblock}
\end{frame}

% FRAME 19 -- Activity 7 solution
\begin{frame}[fragile]{07 | Return four inputs through references -- Solution}
\begin{columns}[T,onlytextwidth]
\begin{column}{0.62\textwidth}
\begin{lstlisting}[style=solutioncpp]
int GetInputData(
    double& S0, double& U, double& D, double& R)
{
    std::cout << "Enter S0 U D R: ";
    std::cin >> S0 >> U >> D >> R;

    if (ValidateInputData(S0, U, D, R) == -1)
    {
        return -1;
    }
    return 0;
}
\end{lstlisting}
\end{column}
\begin{column}{0.35\textwidth}
\small
\explain{Why references}{One function call updates four variables owned by the caller.}
\explain{Meaningful status}{\texttt{0} means success; \texttt{-1} means invalid model data.}
\explain{Caller responsibility}{Check the status before attempting any pricing.}
\end{column}
\end{columns}
\end{frame}

% FRAME 20 -- checkpoint
\begin{frame}{Each Lecture 2 function enters CRR at a specific moment}
\begin{columns}[c,onlytextwidth]
\begin{column}{0.47\textwidth}
\begin{definitionblock}{Before pricing}
\texttt{GetInputData} reads the model inputs.\par
\texttt{ValidateInputData} accepts or rejects them.\par
\texttt{OptionPricer01.cpp} calls \texttt{PriceByCRR} only after valid data.
\end{definitionblock}
\end{column}
\begin{column}{0.49\textwidth}
\begin{definitionblock}{Inside \texttt{PriceByCRR}}
\texttt{RiskNeutProb} supplies $q$. At expiry, \texttt{CalculateAssetPrice} supplies $S(N,i)$ and \texttt{CallPayoff} supplies $H(N,i)$. Backward induction then combines child values.
\end{definitionblock}
\end{column}
\end{columns}
\end{frame}

% FRAME 21 -- Activity 8 problem
\begin{frame}[fragile]{08 | Implement a European call payoff -- Problem}
\problemmeta{Conditional rule}{Complete the function}
\[
h(z)=(z-K)^+=\max(z-K,0).
\]
\begin{lstlisting}[style=solutioncpp]
double CallPayoff(double z, double K)
{
    if (/* in the money */)
    {
        return /* positive payoff */;
    }
    return /* otherwise */;
}
\end{lstlisting}
\begin{definitionblock}{Acceptance cases for $K=100$}
$z=80\to0$, \qquad $z=100\to0$, \qquad $z=120\to20$.
\end{definitionblock}
{\small\color{BootcampBlueDark}\textbf{Boundary question:} Should the condition use \texttt{>} or \texttt{>=}?}
\end{frame}

% FRAME 22 -- Activity 8 solution
\begin{frame}[fragile]{08 | Implement a European call payoff -- Solution}
\begin{columns}[T,onlytextwidth]
\begin{column}{0.57\textwidth}
\begin{lstlisting}[style=solutioncpp]
double CallPayoff(double z, double K)
{
    if (z > K)
    {
        return z - K;
    }
    return 0.0;
}
\end{lstlisting}
\end{column}
\begin{column}{0.40\textwidth}
\small
\explain{At the strike}{$z-K=0$, so both \texttt{>} and \texttt{>=} return zero. \texttt{>} states ``positive payoff'' directly.}
\explain{One responsibility}{The function knows the payoff rule, not how stock prices are generated.}
\explain{Easy to test}{No array or loop is required to verify the three cases.}
\end{column}
\end{columns}
\end{frame}

% FRAME 23 -- Activity 9 problem
\begin{frame}[fragile]{09 | At expiry CRR calls two functions -- Problem}
\problemmeta{Loop + two function calls}{Complete the row}
At expiry $N$, \texttt{PriceByCRR} first calls \texttt{CalculateAssetPrice} for $S(N,i)$, then calls \texttt{CallPayoff} to obtain $H(N,i)=h(S(N,i))$.
\begin{lstlisting}
double Price[SIZE] = { 0.0 };

for (int i = /* first node */; i /* last node */ N; ++i)
{
    double stock = /* compute S(N,i) */;
    Price[i] = /* apply the payoff */;
}
\end{lstlisting}
\begin{definitionblock}{Trace this case}
$S_0=106$, $U=1.15125$, $D=0.86862$, $K=100$, $N=2$.\par
Complete the stock row $[\_\_,\_\_,\_\_]$ and payoff row $[\_\_,\_\_,\_\_]$.
\end{definitionblock}
\end{frame}

% FRAME 24 -- Activity 9 solution
\begin{frame}[fragile]{09 | At expiry CRR calls two functions -- Solution}
\begin{columns}[T,onlytextwidth]
\begin{column}{0.62\textwidth}
\begin{lstlisting}[style=solutioncpp]
double Price[SIZE] = { 0.0 };

for (int i = 0; i <= N; ++i)
{
    double stock = CalculateAssetPrice(
        S0, U, D, N, i);
    Price[i] = CallPayoff(stock, K);
}
\end{lstlisting}
\end{column}
\begin{column}{0.35\textwidth}
\small
\explain{Stock row}{$[79.98,106.00,140.49]$ for $i=0,1,2$.}
\explain{Payoff row}{$[0,6.00,40.49]$ for a call struck at $100$.}
\explain{Invariant}{After the loop, indices $0$ through $N$ hold the complete expiry row.}
\end{column}
\end{columns}
\end{frame}

% FRAME 25 -- trace
\begin{frame}{After expiry, CRR needs $q$ and two child values}
\texttt{PriceByCRR} calls \texttt{RiskNeutProb(U, D, R)} once; it returns $q\approx0.48413$. With $R=1.00545$, begin with the terminal row $[0,6.00,40.49]$.
\begin{center}
\renewcommand{\arraystretch}{1.35}
\begin{tabular}{ccl}
\toprule
Step & Active values & Update rule \\
\midrule
$n=2$ & $[0,6.00,40.49]$ & terminal call payoffs \\
$n=1$ & $[2.8890,22.5745]$ & $H(1,i)=\frac{qH(2,i+1)+(1-q)H(2,i)}{R}$ \\
$n=0$ & $[12.3521]$ & same local rule \\
\bottomrule
\end{tabular}
\end{center}
\begin{propositionblock}{Why one array is enough}
After computing row $n$, the old row $n+1$ is never needed again. We overwrite it from left to right.
\end{propositionblock}
\end{frame}

% FRAME 26 -- Activity 10 problem
\begin{frame}[fragile]{10 | Complete backward induction -- Problem}
\problemmeta{Two loops + in-place update}{Complete four blanks}
\begin{lstlisting}
for (int n = /* first backward step */;
     n /* continue through root */ 0; --n)
{
    for (int i = 0; i /* nodes on row n */ n; ++i)
    {
        Price[i] =
            (q * /* up child */
             + (1.0 - q) * /* down child */) / R;
    }
}
return Price[0];
\end{lstlisting}
\begin{definitionblock}{Index map}
At node $(n,i)$, the down child is \texttt{Price[i]} and the up child is \texttt{Price[i+1]}.
\end{definitionblock}
\end{frame}

% FRAME 27 -- Activity 10 solution
\begin{frame}[fragile]{10 | Complete backward induction -- Solution}
\begin{columns}[T,onlytextwidth]
\begin{column}{0.61\textwidth}
\begin{lstlisting}[style=solutioncpp]
for (int n = N - 1; n >= 0; --n)
{
    for (int i = 0; i <= n; ++i)
    {
        Price[i] =
            (q * Price[i + 1]
             + (1.0 - q) * Price[i]) / R;
    }
}
return Price[0];
\end{lstlisting}
\end{column}
\begin{column}{0.36\textwidth}
\small
\explain{Outer loop}{Expiry is already known, so begin at $N-1$ and finish at the root.}
\explain{Inner loop}{Row $n$ contains indices $0$ through $n$.}
\explain{Final state}{After $n=0$, \texttt{Price[0]} is $H(0,0)$.}
\end{column}
\end{columns}
\end{frame}

% FRAME 28 -- Activity 11 problem
\begin{frame}[fragile]{11 | Find the overwrite bug -- Problem}
\problemmeta{Data dependency}{The formulas are right; the result is wrong}
\begin{lstlisting}
for (int n = N - 1; n >= 0; --n)
{
    for (int i = n; i >= 0; --i)
    {
        Price[i] =
            (q * Price[i + 1]
             + (1.0 - q) * Price[i]) / R;
    }
}
\end{lstlisting}
\begin{definitionblock}{Trace only row $n=1$}
Start from \texttt{Price = [0, 6.00, 40.49]}. Which entry is overwritten first? Which later calculation reads that changed entry as if it were still a child?
\end{definitionblock}
{\small\color{BootcampBlueDark}\textbf{Repair:} Change only the inner-loop header.}
\end{frame}

% FRAME 29 -- Activity 11 solution
\begin{frame}[fragile]{11 | Find the overwrite bug -- Solution}
\begin{columns}[T,onlytextwidth]
\begin{column}{0.56\textwidth}
\begin{lstlisting}[style=solutioncpp]
for (int n = N - 1; n >= 0; --n)
{
    for (int i = 0; i <= n; ++i)
    {
        Price[i] =
            (q * Price[i + 1]
             + (1.0 - q) * Price[i]) / R;
    }
}
\end{lstlisting}
\end{column}
\begin{column}{0.41\textwidth}
\small
\explain{What went wrong}{With decreasing \texttt{i}, \texttt{Price[1]} changes before node $(1,0)$ reads it as its up child.}
\explain{Why increasing works}{Updating \texttt{Price[0]} cannot damage \texttt{Price[1]}, which is still needed next.}
\explain{General lesson}{When overwriting an array, loop direction is part of correctness.}
\end{column}
\end{columns}
\end{frame}

% FRAME 30 -- Activity 12 problem A
\begin{frame}[fragile]{12 | PriceByCRR begins with validation and $q$ -- Problem A}
\problemmeta{Guard + initialization}{Complete the preconditions}
\begin{lstlisting}[style=tinycpp]
double PriceByCRR(double S0, double U, double D,
                  double R, const int N, double K)
{
    const int SIZE = 9;
    if (N < 1 || N >= SIZE || S0 <= 0.0 || K < 0.0
        || /* require 0 < D < R < U */)
    {
        std::cerr << "Invalid model data" << std::endl;
        return -1.0;
    }

    double q = /* call RiskNeutProb(U, D, R) */;
    double Price[SIZE] = { 0.0 };
    // Next: construct the terminal row.
\end{lstlisting}
\begin{definitionblock}{Design check}
Why does \texttt{SIZE = 9} allow $N=8$ but reject $N=9$?
\end{definitionblock}
\end{frame}

% FRAME 31 -- Activity 12 problem B
\begin{frame}[fragile]{12 | PriceByCRR builds expiry then moves backward -- Problem B}
\problemmeta{Terminal row + backward pass}{Connect the completed pieces}
\begin{lstlisting}[style=tinycpp]
    for (int i = 0; i <= N; ++i)
    {
        Price[i] = /* payoff at S(N,i) */;
    }

    for (int n = /* ? */; n >= 0; --n)
    {
        for (int i = 0; i <= n; ++i)
        {
            Price[i] = /* one CRR update */;
        }
    }
    return /* root value */;
}
\end{lstlisting}
\begin{definitionblock}{Before the reveal}
Point to the base case, the recurrence, and the final answer in this function.
\end{definitionblock}
\end{frame}

% FRAME 32 -- Activity 12 solution A
\begin{frame}[fragile]{12 | PriceByCRR begins with validation and $q$ -- Solution A}
\begin{lstlisting}[style=tinycpp]
double PriceByCRR(double S0, double U, double D,
                  double R, const int N, double K)
{
    const int SIZE = 9;
    if (N < 1 || N >= SIZE || S0 <= 0.0 || K < 0.0
        || D <= 0.0 || U <= D || R <= D || R >= U)
    {
        std::cerr << "Invalid model data" << std::endl;
        return -1.0;
    }

    double q = RiskNeutProb(U, D, R);
    double Price[SIZE] = { 0.0 };
    for (int i = 0; i <= N; ++i)
    {
        Price[i] = CallPayoff(
            CalculateAssetPrice(S0, U, D, N, i), K);
    }
\end{lstlisting}
{\small The validation establishes safe ranges; the terminal loop establishes the base case.}
\end{frame}

% FRAME 33 -- Activity 12 solution B
\begin{frame}[fragile]{12 | PriceByCRR builds expiry then moves backward -- Solution B}
\begin{lstlisting}[style=tinycpp]
    for (int n = N - 1; n >= 0; --n)
    {
        for (int i = 0; i <= n; ++i)
        {
            Price[i] =
                (q * Price[i + 1]
                 + (1.0 - q) * Price[i]) / R;
        }
    }
    return Price[0];
}
\end{lstlisting}
\begin{columns}[T,onlytextwidth]
\begin{column}{0.48\textwidth}
\begin{propositionblock}{Shared example}
$N=8$ returns approximately $21.68$ for $S_0=106$, $U=1.15125$, $D=0.86862$, $R=1.00545$, $K=100$.
\end{propositionblock}
\end{column}
\begin{column}{0.48\textwidth}
\begin{illustrationblock}{Explain the memory}
At any moment, the active prefix of \texttt{Price} represents one complete time level.
\end{illustrationblock}
\end{column}
\end{columns}
\end{frame}

% FRAME 34 -- delta hedge concept
\begin{frame}{A CRR node gives both an option value and a delta hedge}
At a node $(n,i)$, replicate the option over one step with shares of stock and risk-free cash:
\[
\Delta(n,i)=
\frac{H(n+1,i+1)-H(n+1,i)}
     {S(n+1,i+1)-S(n+1,i)}.
\]
\begin{columns}[T,onlytextwidth]
\begin{column}{0.47\textwidth}
\begin{definitionblock}{Interpretation}
Hold $\Delta$ shares. Put the remaining amount in the risk-free account (or borrow it when it is negative). The portfolio has the same up and down values as the option.
\end{definitionblock}
\end{column}
\begin{column}{0.49\textwidth}
\begin{definitionblock}{Where the stock values come from}
\texttt{CalculateAssetPrice} provides the two child prices $S(n+1,i)$ and $S(n+1,i+1)$. The array \texttt{Price} holds the corresponding option values.
\end{definitionblock}
\end{column}
\end{columns}
\end{frame}

% FRAME 35 -- Activity 13 delta implementation
\begin{frame}[fragile]{13 | Compute delta before overwriting \texttt{Price[i]}}
\begin{columns}[T,onlytextwidth]
\begin{column}{0.59\textwidth}
\begin{lstlisting}[style=solutioncpp]
double Su = CalculateAssetPrice(S0, U, D, n + 1, i + 1);
double Sd = CalculateAssetPrice(S0, U, D, n + 1, i);
double delta = (Price[i + 1] - Price[i]) / (Su - Sd);

Price[i] = (q * Price[i + 1]
            + (1.0 - q) * Price[i]) / R;
\end{lstlisting}
\end{column}
\begin{column}{0.38\textwidth}
\small
\explain{Lecture 2 $N=2$ root}{$S_d=92.0737$, $S_u=122.0325$, $H_d=2.8890$, $H_u=22.5745$.}
\explain{Hedge}{$\Delta\approx0.6571$: hold $0.6571$ share. Cash is $H(0,0)-\Delta S_0\approx-57.30$, so borrow $57.30$.}
\explain{Order matters}{Calculate delta before replacing \texttt{Price[i]}; it needs both child option values.}
\end{column}
\end{columns}
\end{frame}

% FRAME 36 -- Lecture 2 driver program
\begin{frame}[fragile]{OptionPricer01.cpp runs the Lecture 2 example}
\begin{columns}[T,onlytextwidth]
\begin{column}{0.67\textwidth}
\begin{lstlisting}[style=tinycpp]
#include "BinomialTreeModel.h"
#include "Option01.h"
#include <iostream>
#include <iomanip>

int main()
{
    double u = 1.15125, d = 0.86862;
    double r = 1.00545, s0 = 106.00;
    double k = 100.00;
    const int N = 8;
    double optionPrice =
        fre::PriceByCRR(s0, u, d, r, N, k);
    std::cout << "European call option price = "
              << std::fixed << std::setprecision(2)
              << optionPrice << std::endl;
    return 0;
}
\end{lstlisting}
\end{column}
\begin{column}{0.30\textwidth}
\small
\begin{propositionblock}{Expected output}
\texttt{European call option price = 21.68}
\end{propositionblock}
\explain{Same file roles}{The model helpers live in \texttt{BinomialTreeModel}; the payoff and pricer live in \texttt{Option01}.}
\end{column}
\end{columns}
\end{frame}

% FRAME 37 -- tests
\begin{frame}{Tests should target boundaries and invariants}
\scriptsize
\begin{center}
\renewcommand{\arraystretch}{1.18}
\begin{tabular}{@{}p{1.55cm}p{4.15cm}p{2.55cm}p{2.75cm}@{}}
\toprule
Case & Change & Expected observation & Bug targeted \\
\midrule
Lecture example & $S_0=106,U=1.15125,D=.86862,R=1.00545,K=100,N=8$ & $21.68$ & full data flow \\
At strike & call payoff $z=K=100$ & $0$ & payoff boundary \\
No-arbitrage edge & $R=D$ & rejected & wrong inequality \\
Capacity edge & $N=9$ with \texttt{SIZE=9} & rejected & out-of-bounds write \\
Zero strike & set $K=0$ & equal to $S_0=106$ & discount/update error \\
\bottomrule
\end{tabular}
\end{center}
\begin{propositionblock}{A useful explanation has three parts}
{\small State the input, predict the result or invariant, then identify the line that would be suspect if the test failed.}
\end{propositionblock}
\end{frame}

% FRAME 38 -- core close
\begin{frame}{One algorithm now has a line-by-line explanation}
\begin{columns}[c,onlytextwidth]
\begin{column}{0.46\textwidth}
\centering
Validate inputs\par
$\downarrow$\par
Compute $q$\par
$\downarrow$\par
Build terminal payoffs\par
$\downarrow$\par
Move backward in place\par
$\downarrow$\par
Return \texttt{Price[0]}
\end{column}
\begin{column}{0.50\textwidth}
\begin{propositionblock}{Exit ticket}
\begin{enumerate}
\item Why must $R$ lie strictly between $D$ and $U$?
\item Why does the inner CRR loop move from left to right?
\item Which function would you test first after changing the payoff from call to put?
\end{enumerate}
\end{propositionblock}
\medskip
If you can answer all three, you can debug the pricer rather than merely rerun it.
\end{column}
\end{columns}
\end{frame}

% RESERVE PRACTICE -- adapted directly from Recitation 1 Deck 5.
\begin{frame}{Additional practice reuses the same coding habits}
\begin{columns}[c,onlytextwidth]
\begin{column}{0.46\textwidth}
\begin{definitionblock}{For each problem}
\begin{enumerate}
\item Trace the example by hand.
\item Complete the function.
\item Test one boundary case.
\item Explain the loop invariant.
\end{enumerate}
\end{definitionblock}
\end{column}
\begin{column}{0.50\textwidth}
\begin{illustrationblock}{Choose a useful next challenge}
Count and accumulate first, then move to nested loops or a one-pass financial algorithm.
\end{illustrationblock}
\medskip
These problems come from Recitation 1 Deck 5 and can be used in class or after the core CRR lab.
\smallskip\par
{\small Use \texttt{<vector>}; the last two solutions also use \texttt{<algorithm>}.}
\end{column}
\end{columns}
\end{frame}

\begin{frame}[fragile]{Practice 01 | Count positive values -- Problem}
\problemmeta{Counter + condition}{Complete the function}
Return how many elements of an integer vector are strictly positive. Zero does not count.
\begin{definitionblock}{Examples}
\texttt{[-2, 0, 5, 3, -1]} $\to$ \texttt{2}\qquad
\texttt{[0, -4]} $\to$ \texttt{0}
\end{definitionblock}
\begin{lstlisting}[style=tinycpp]
int countPositive(const std::vector<int>& nums)
{
    int count = /* ? */;
    for (int i = 0; i < static_cast<int>(nums.size()); ++i)
    {
        if (/* current value is positive */)
        {
            /* update count */;
        }
    }
    return /* ? */;
}
\end{lstlisting}
\end{frame}

\begin{frame}[fragile]{Practice 01 | Count positive values -- Solution}
\begin{columns}[T,onlytextwidth]
\begin{column}{0.61\textwidth}
\begin{lstlisting}[style=solutioncpp]
int countPositive(const std::vector<int>& nums)
{
    int count = 0;
    int n = static_cast<int>(nums.size());
    for (int i = 0; i < n; ++i)
    {
        if (nums[i] > 0)
        {
            ++count;
        }
    }
    return count;
}
\end{lstlisting}
\end{column}
\begin{column}{0.36\textwidth}
\small
\explain{Invariant}{\texttt{count} equals the number of positive values already visited.}
\explain{Boundary}{An empty vector returns zero because the loop never runs.}
\explain{Cost}{$O(n)$ time and $O(1)$ extra space.}
\end{column}
\end{columns}
\end{frame}

\begin{frame}[fragile]{Practice 02 | Running sum -- Problem}
\problemmeta{Accumulate a result}{Complete the function}
Return a vector whose value at index $i$ is the sum of inputs from index $0$ through $i$.
\begin{definitionblock}{Example}
\texttt{[2, 1, 3, 4]} $\to$ \texttt{[2, 3, 6, 10]}
\end{definitionblock}
\begin{lstlisting}[style=tinycpp]
std::vector<int> runningSum(const std::vector<int>& nums)
{
    int n = static_cast<int>(nums.size());
    std::vector<int> result(n);
    int total = /* ? */;
    for (int i = 0; i < n; ++i)
    {
        /* update total */;
        result[i] = /* ? */;
    }
    return result;
}
\end{lstlisting}
\end{frame}

\begin{frame}[fragile]{Practice 02 | Running sum -- Solution}
\begin{columns}[T,onlytextwidth]
\begin{column}{0.61\textwidth}
\begin{lstlisting}[style=solutioncpp]
std::vector<int> runningSum(
    const std::vector<int>& nums)
{
    int n = static_cast<int>(nums.size());
    std::vector<int> result(n);
    int total = 0;
    for (int i = 0; i < n; ++i)
    {
        total += nums[i];
        result[i] = total;
    }
    return result;
}
\end{lstlisting}
\end{column}
\begin{column}{0.36\textwidth}
\small
\explain{Invariant}{After index $i$, \texttt{total} equals \texttt{nums[0]+...+nums[i]}.}
\explain{Reuse}{The next iteration adds one value instead of recomputing the prefix.}
\explain{Cost}{$O(n)$ time and $O(n)$ output space.}
\end{column}
\end{columns}
\end{frame}

\begin{frame}[fragile]{Practice 03 | Richest customer -- Problem}
\problemmeta{Nested loops + maximum}{Complete the resets and updates}
Each row contains one customer's balances. Return the largest row total.
\smallskip\par
\textbf{Example:} \texttt{[[2, 3], [1, 8], [4, 4]]} $\to$ \texttt{9}
\begin{lstlisting}[style=tinycpp]
int maximumWealth(
    const std::vector<std::vector<int>>& accounts)
{
    int best = 0;
    for (int r = 0; r < static_cast<int>(accounts.size()); ++r)
    {
        int total = /* reset for this row */;
        for (int c = 0; c < static_cast<int>(accounts[r].size()); ++c)
        {
            /* add one balance */;
        }
        best = /* keep the larger value */;
    }
    return best;
}
\end{lstlisting}
\end{frame}

\begin{frame}[fragile]{Practice 03 | Richest customer -- Solution}
\begin{columns}[T,onlytextwidth]
\begin{column}{0.63\textwidth}
\begin{lstlisting}[style=tinycpp]
int maximumWealth(
    const std::vector<std::vector<int>>& accounts)
{
    int best = 0;
    int rows = static_cast<int>(accounts.size());
    for (int r = 0; r < rows; ++r)
    {
        int total = 0;
        int banks = static_cast<int>(accounts[r].size());
        for (int c = 0; c < banks; ++c)
        {
            total += accounts[r][c];
        }
        best = std::max(best, total);
    }
    return best;
}
\end{lstlisting}
\end{column}
\begin{column}{0.34\textwidth}
\small
\explain{Inner loop}{Computes one customer's total.}
\explain{Outer loop}{Compares that total with the best completed row.}
\explain{Critical reset}{\texttt{total = 0} belongs inside the outer loop.}
\end{column}
\end{columns}
\end{frame}

\begin{frame}[fragile]{Practice 04 | Best time to trade -- Problem}
\problemmeta{Track the best value so far}{Complete one pass}
Return the maximum profit from one buy followed by one later sell. Return zero when no profitable trade exists.
{\small Assume the vector contains at least one price.\par}
\smallskip\par
\textbf{Examples:} \texttt{[7,1,5,3,6,4]} $\to$ \texttt{5}\qquad
\texttt{[7,6,4,3,1]} $\to$ \texttt{0}
\begin{lstlisting}[style=tinycpp]
int maxProfit(const std::vector<int>& prices)
{
    int minPrice = prices[0];
    int best = 0;
    for (int i = 1; i < static_cast<int>(prices.size()); ++i)
    {
        int profit = /* sell today */;
        best = /* best profit so far */;
        minPrice = /* cheapest price so far */;
    }
    return best;
}
\end{lstlisting}
\end{frame}

\begin{frame}[fragile]{Practice 04 | Best time to trade -- Solution}
\begin{columns}[T,onlytextwidth]
\begin{column}{0.61\textwidth}
\begin{lstlisting}[style=solutioncpp]
int maxProfit(const std::vector<int>& prices)
{
    int n = static_cast<int>(prices.size());
    int minPrice = prices[0];
    int best = 0;
    for (int i = 1; i < n; ++i)
    {
        int profit = prices[i] - minPrice;
        best = std::max(best, profit);
        minPrice = std::min(minPrice, prices[i]);
    }
    return best;
}
\end{lstlisting}
\end{column}
\begin{column}{0.36\textwidth}
\small
\explain{Invariant}{Before day $i$, \texttt{minPrice} is the cheapest earlier price.}
\explain{Order matters}{Evaluate today's sale before adding today's price to the past.}
\explain{Cost}{$O(n)$ time and $O(1)$ extra space.}
\end{column}
\end{columns}
\end{frame}

\begin{frame}{The same habits transfer beyond CRR}
\begin{columns}[T,onlytextwidth]
\begin{column}{0.47\textwidth}
\begin{definitionblock}{State a loop invariant}
Name what each accumulator, index, or active array prefix means before and after an iteration.
\end{definitionblock}
\begin{definitionblock}{Protect boundaries}
Test empty inputs, equality cases, first and last indices, and invalid model parameters.
\end{definitionblock}
\end{column}
\begin{column}{0.49\textwidth}
\begin{propositionblock}{Explain before optimizing}
A correct explanation connects the specification, the stored state, the update, and a test that could falsify the code.
\end{propositionblock}
\medskip
That is the common skill behind every exercise in this deck.
\end{column}
\end{columns}
\end{frame}

\end{document}
